% Generated draft: requires mathematical and source review.
\documentclass[11pt]{article}
\usepackage[margin=25mm]{geometry}
\usepackage{fontspec}
\setmainfont{Latin Modern Roman}
\usepackage{amsmath,amssymb,mathtools,unicode-math}
\setmathfont{Latin Modern Math}
\usepackage{enumitem,needspace,xurl,hyperref}
\setlength{\parindent}{0pt}
\setlength{\parskip}{5pt}
\setlength{\emergencystretch}{3em}
\newcommand{\N}{\mathbb{N}}
\newcommand{\Z}{\mathbb{Z}}
\newcommand{\Q}{\mathbb{Q}}
\newcommand{\R}{\mathbb{R}}
\newcommand{\C}{\mathbb{C}}
\newcommand{\F}{\mathbb{F}}
\newcommand{\A}{\mathbb{A}}
\newcommand{\PP}{\mathbb{P}}
\newcommand{\E}{\mathbb{E}}
\newcommand{\eps}{\varepsilon}
\newcommand{\dd}{\,\mathrm{d}}
\begin{document}
\section*{UnsolvedMath ID: 30002402}
\textbf{Source problem code:} catalogue entry
\subsection{Definitions and mathematical statement}
The supplied question concerns compact Calabi--Yau threefolds and their corresponding Pfaffian, determinantal, or hybrid Landau--Ginzburg phases arising from non-abelian linear sigma models. The record does not specify the models, their parameters, the phase correspondence, or the mathematical meaning of their D-brane categories.

To express the intended assertion schematically, let \(\mathcal M\) denote the collection of non-abelian linear sigma models intended by the source. For \(M\in\mathcal M\), let \(X\) be a compact Calabi--Yau threefold phase of \(M\), and let \(P\) be a corresponding Pfaffian, determinantal, or hybrid Landau--Ginzburg phase. Write \(\mathcal D_M(X)\) and \(\mathcal D_M(P)\) for the intended D-brane categories; these symbols are placeholders, since the record supplies no categorical definitions.

\emph{Question, in schematic form.} For every such model \(M\) and every intended corresponding pair \((X,P)\), does there exist an equivalence
\[
\mathcal D_M(X)\simeq\mathcal D_M(P)?
\]
The collection \(\mathcal M\), the admissible parameters and phase pairs, and the meaning of equivalence must be specified before this schema becomes a fully determined mathematical conjecture. In particular, the supplied record does not identify the D-brane categories with derived categories of coherent sheaves or with categories of matrix factorizations, and no such identification is assumed here. It also does not specify the base field, smoothness conventions, treatment of singular phases, or any grading, equivariance, or enhancement required for the categories.
\subsection{Short English statement}
The question asks whether corresponding compact Calabi--Yau threefold and Pfaffian, determinantal, or hybrid Landau--Ginzburg phases of the intended non-abelian linear sigma models have equivalent D-brane categories. The models and categories are not specified sufficiently to produce a unique precise conjecture. The supplied partially solved label and claims about selected examples require independent review.
\subsection{Sources}
\begin{itemize}
\item Local UnsolvedMath registry, supplied record with permanent ID \(30002402\), source problem code \texttt{catalogue entry}, title ``D-Brane Equivalence Across Calabi--Yau Phases,'' in the Oberwolfach Reports Open Problems collection. The project identifies \texttt{data/unsolvedmath/problems.json} and the supplementary \texttt{catalog/problems.json} as registry locations; the supplied extract does not identify which contains this record. No individual proposer is supplied; the proposed year is recorded as 2013.
\item Supplied source citation: \emph{Matrix Factorizations in Algebra, Geometry, and Physics} (2014), Oberwolfach Reports, pp.~2501--2552, DOI \url{https://doi.org/10.4171/owr/2013/44}. The background attributes the problem to mathematical workshop and open-problem discussions documented in Oberwolfach Reports.
\item Supplied literature source: ``Oberwolfach Report 44,'' \url{https://publications.mfo.de/bitstream/handle/mfo/3375/OWR_2013_44.pdf?sequence=1}.
\item Supplied literature source: ``Hori--Knapp,'' \url{https://arxiv.org/abs/1308.6265}.
\item Supplied literature source: ``Grade restriction rules,'' \url{https://doi.org/10.1007/s11401-017-1103-8}.
\item Supplied dated literature curation, checked \texttt{2026-08-21}: status ``partially solved.'' Its assessment states that the phase construction predicts equivalent D-brane categories, that later work proves selected examples, and that the broad conjecture for the stated exotic Calabi--Yau families is not proved. This is supplied dated curation derived from the OWR material, not an independently verified literature assessment; its status and cited claims require independent verification.
\item The registry labels the statement ``exact'' and reports that its mathematical content was checked against the official Oberwolfach report, with notation, context, and question form normalized without changing the source-level claim. This verification claim is supplied metadata and has not been independently checked for this draft.
\end{itemize}
\end{document}
